% ============================================================================
% LANGENTWURF: Wiederholung + Stundenleistung - Mi casa
% ============================================================================
% Sequenz 3, Block d: Wiederholung + Stundenleistung
% Fachfarbe: #c41e3a (Karminrot)
% ============================================================================

\documentclass[11pt,a4paper]{article}

% ============================================================================
% PACKAGES
% ============================================================================
\usepackage[utf8]{inputenc}
\usepackage[T1]{fontenc}
\usepackage[ngerman]{babel}
\usepackage{geometry}
\usepackage{fancyhdr}
\usepackage{lastpage}
\usepackage{xcolor}
\usepackage{tcolorbox}
\usepackage{enumitem}
\usepackage{tabularx}
\usepackage{longtable}
\usepackage{array}
\usepackage{booktabs}
\usepackage{hyperref}
\usepackage{ifthen}
\usepackage{amssymb}

% ============================================================================
% KONFIGURATION
% ============================================================================

\newcommand{\plantier}{1}
\newcommand{\fachid}{spanisch}
\newcommand{\fach}{Spanisch}
\newcommand{\fachfarbe}{c41e3a}

% ============================================================================
% VARIABLEN
% ============================================================================

\newcommand{\jahrgangsstufe}{7}
\newcommand{\schulart}{Gesamtschule}
\newcommand{\stundenthema}{Wiederholung + Stundenleistung}
\newcommand{\reihenthema}{Mi casa}
\newcommand{\stundennummer}{3d}
\newcommand{\reihenstunden}{4}
\newcommand{\unterrichtsdatum}{}
\newcommand{\unterrichtszeit}{}
\newcommand{\raum}{}

\newcommand{\lehrkraft}{N.N.}
\newcommand{\ausbildungsstand}{Lehrkraft}

\newcommand{\lerngruppengroesse}{24}
\newcommand{\maennlich}{12}
\newcommand{\weiblich}{12}
\newcommand{\divers}{0}

\newcommand{\gerniveau}{A1}
\newcommand{\lehrwerk}{Que pasa 1}
\newcommand{\lehrwerkseite}{S. 36-45}

% ============================================================================
% LAYOUT
% ============================================================================
\geometry{left=2.5cm, right=2.5cm, top=2.5cm, bottom=2.5cm}
\definecolor{fach}{HTML}{\fachfarbe}
\definecolor{gray}{HTML}{666666}
\definecolor{lightgray}{HTML}{f5f5f5}
\definecolor{successgreen}{HTML}{2a7a4b}
\definecolor{warningorange}{HTML}{b35c00}

\tcbuselibrary{skins,breakable}

\newtcolorbox{infobox}[1][]{
    colback=lightgray,
    colframe=fach,
    fonttitle=\bfseries,
    title=#1,
    breakable
}

\newtcolorbox{scorebox}{
    colback=white,
    colframe=fach,
    fonttitle=\bfseries,
    title=Didaktik-Score,
    breakable
}

\pagestyle{fancy}
\fancyhf{}
\fancyhead[L]{\small\textcolor{gray}{\fach{} -- Jg. \jahrgangsstufe{} -- \stundenthema}}
\fancyhead[R]{\small\textcolor{gray}{Langentwurf Tier 1}}
\fancyfoot[C]{\small\thepage{} / \pageref{LastPage}}
\renewcommand{\headrulewidth}{0.4pt}
\renewcommand{\footrulewidth}{0pt}

% ============================================================================
% DOKUMENT
% ============================================================================
\begin{document}

% ============================================================================
% DECKBLATT
% ============================================================================
\begin{titlepage}
    \centering
    \vspace*{2cm}

    {\large\textcolor{gray}{\schulart}}\\[2cm]

    {\Huge\textcolor{fach}{\textbf{Langentwurf}}}\\[0.5cm]
    {\large Tier 1 -- Vollstandig}\\[2cm]

    {\LARGE\textbf{\stundenthema}}\\[0.5cm]
    {\large im Rahmen der Unterrichtsreihe}\\[0.3cm]
    {\Large\textit{\reihenthema}}\\[2cm]

    \begin{tabular}{rl}
        \textbf{Fach:} & \fach{} (\gerniveau) \\[0.3cm]
        \textbf{Klasse:} & \jahrgangsstufe \\[0.3cm]
        \textbf{Lehrwerk:} & \lehrwerk, \lehrwerkseite \\[0.3cm]
        \textbf{Datum:} & \underline{\hspace{3cm}} \\[0.3cm]
        \textbf{Zeit:} & \underline{\hspace{3cm}} (Doppelstunde) \\[0.3cm]
        \textbf{Raum:} & \underline{\hspace{3cm}} \\[0.3cm]
    \end{tabular}

    \vfill

    {\large Lehrkraft: \lehrkraft}\\[0.3cm]
    {\normalsize\textcolor{gray}{\ausbildungsstand}}\\[1cm]

\end{titlepage}

% ============================================================================
% INHALTSVERZEICHNIS
% ============================================================================
\tableofcontents
\newpage

% ============================================================================
% 1. BEDINGUNGSANALYSE
% ============================================================================
\section{Bedingungsanalyse}

\subsection{Lerngruppe}
\begin{tabular}{ll}
    Kursgröße: & \lerngruppengroesse{} SuS (\maennlich{} m / \weiblich{} w / \divers{} d) \\
    Jahrgangsstufe: & \jahrgangsstufe \\
    Sprachniveau: & \gerniveau{} (GER) \\
    Fremdsprache: & 2. Fremdsprache \\
\end{tabular}

\subsection{Differenzierung (intern)}

\textbf{WICHTIG:} Keine Sternchen auf Schulermaterial!

\begin{tabular}{|l|p{4cm}|p{4cm}|p{4cm}|}
\hline
& \textbf{[B] Basis} & \textbf{[S] Standard} & \textbf{[E] Erweiterung} \\
\hline
\textbf{Ziel} & Berufsreife & Mittlere Reife & Gymnasium \\
\hline
\textbf{Inhalt} & Rezeptiv, Zuordnung & Produktion mit Hilfe & Freie Produktion \\
\hline
\textbf{Hilfen} & Wortlisten, Bilder & Teils vorgegeben & Nur Impulse \\
\hline
\end{tabular}

\subsection{Besonderheiten}
\begin{itemize}
    \item \textbf{LRS:} 2 SuS (Nachteilsausgleich: +5 Min. bei Stundenleistung)
    \item \textbf{DaZ:} 1 SuS (Wortschatz deutsch-spanisch vorab besprechen)
    \item \textbf{Leistungsstark:} 3 SuS (Aufgabe 4 erweitert)
\end{itemize}

% ============================================================================
% 2. SACHANALYSE
% ============================================================================
\section{Sachanalyse}

\subsection{Sprachliche Strukturen}
\begin{itemize}
    \item \textbf{Grammatik:} \textit{hay} (es gibt), Adjektivanpassung (Genus/Numerus)
    \item \textbf{Wortschatz:} Raume (la cocina, el salon, el dormitorio, el bano, el jardin), Mobel (la cama, la mesa, la silla, el sofa, el armario), Farben (rojo, azul, verde, amarillo, blanco, negro, gris, marron)
    \item \textbf{Phonetik:} Aussprache von ll, n, j; Betonung bei Adjektiven
\end{itemize}

\subsection{Kontrastive Betrachtung (Deutsch $\leftrightarrow$ Spanisch)}
\begin{itemize}
    \item \textbf{Positive Transfers:} Ahnliche Wortstellung bei Farbadjektiven nach Nomen
    \item \textbf{Interferenzen:} \textit{hay} = ,,es gibt'' (unveranderlich, kein Artikel danach)
    \item \textbf{Falsche Freunde:} \textit{la habitacion} = Zimmer (nicht ,,Habitation'')
\end{itemize}

\subsection{Kulturelle Aspekte}
\begin{itemize}
    \item Typische spanische Hauser (patio, terraza)
    \item Wohnverhaltnisse in Spanien vs. Deutschland
\end{itemize}

% ============================================================================
% 3. DIDAKTISCHE ANALYSE
% ============================================================================
\section{Didaktische Analyse}

\subsection{Stellung in der Sequenz}
Dies ist die \textbf{4. Doppelstunde} der Sequenz ,,\reihenthema'' (4 Doppelstunden gesamt).

\textbf{Vorangegangene Stunden:}
\begin{itemize}
    \item 03a: Raume im Haus (habitaciones)
    \item 03b: \textit{hay} + Mobel (muebles)
    \item 03c: Farben + Adjektivanpassung (colores)
\end{itemize}

\subsection{Kompetenzziele (Rahmenplan MV)}
\begin{itemize}
    \item \textbf{Kommunikativ:} Leseverstehen, Schreiben (beschreiben)
    \item \textbf{Sprachlich:} \textit{hay}, Adjektivkongruenz, Wortschatz Haus
    \item \textbf{Interkulturell:} Wohnen in Spanien
    \item \textbf{Methodisch:} Selbststandige Textproduktion
\end{itemize}

\subsection{Legitimation}
\textbf{Rahmenplan Spanisch MV (Kl. 7-10):}
\begin{itemize}
    \item Kompetenzbereich: Schreiben -- einfache Texte zu vertrauten Themen
    \item Standard: Alltagssituationen beschreiben (Wohnung, Zimmer)
\end{itemize}

% ============================================================================
% 4. STUNDENZIELE
% ============================================================================
\section{Stundenziele}

\subsection{Grobziel}
Die SuS festigen ihre Kenntnisse zu Raumvokabular, \textit{hay}, Mobeln und Farben und weisen diese in einer Stundenleistung nach.

\subsection{Feinziele (operationalisiert)}
\begin{tabular}{|c|p{8cm}|c|c|}
\hline
\textbf{Nr.} & \textbf{Die SuS konnen...} & \textbf{AFB} & \textbf{Diff.} \\
\hline
FZ1 & Raumvokabular korrekt zuordnen & I & alle \\
\hline
FZ2 & Satze mit \textit{hay} + Mobel bilden & II & alle \\
\hline
FZ3 & Farbadjektive an Genus/Numerus anpassen & II & alle \\
\hline
FZ4 & Ein Zimmer eigenstandig beschreiben & II/III & [S]/[E] \\
\hline
\end{tabular}

% ============================================================================
% 5. METHODISCHE ANALYSE
% ============================================================================
\section{Methodische Analyse}

\subsection{Phasierung (Doppelstunde = 2 x 45 Min.)}
\begin{enumerate}
    \item \textbf{1. Stunde (45 Min.):} Wiederholung
    \begin{itemize}
        \item Einstieg: Schnelles Vokabelquiz (mundlich)
        \item Wiederholung: Partnerarbeit mit Lernpfad/AB
        \item Ubung: Gemischte Aufgaben zu allen Themen
        \item Sicherung: Klaren von Fragen
    \end{itemize}
    \item \textbf{2. Stunde (45 Min.):} Stundenleistung
    \begin{itemize}
        \item Vorbereitung: Materialausgabe, Ruhe
        \item Durchfuhrung: Stundenleistung (25 Min.)
        \item Abschluss: Einsammeln, Ausblick
    \end{itemize}
\end{enumerate}

\subsection{Medien}
\begin{itemize}
    \item Tafel/Whiteboard
    \item Lehrwerk: \lehrwerk, \lehrwerkseite
    \item Wiederholungs-AB: AB-03d.pdf
    \item Lernpfad: LP-03d.html (optional, digital)
    \item Stundenleistung: SL-03d.pdf
\end{itemize}

% ============================================================================
% 6. VERLAUFSPLANUNG
% ============================================================================
\section{Verlaufsplanung}

\textbf{1. Stunde: Wiederholung (45 Min.)}

\begin{longtable}{|p{1cm}|p{1.5cm}|p{4cm}|p{3cm}|p{2cm}|p{2cm}|}
\hline
\textbf{Zeit} & \textbf{Phase} & \textbf{Lehrerhandeln} & \textbf{Schulerhandeln} & \textbf{SF} & \textbf{Medien} \\
\hline
\endhead

0--5 & Einstieg & Begrußung, Ankun\-digung: ,,Heute wiederholen wir alles zu mi casa'' & Zuhoren & Plenum & Tafel \\
\hline

5--10 & Aktivie\-rung & Schnelles Vokabelquiz: Raume zeigen, SuS benennen & Antworten auf Spanisch & Plenum & Bilder \\
\hline

10--25 & Wieder\-holung & Lernpfad/AB-Bearbeitung anleiten, unterstutzen & Aufgaben bearbeiten & EA/PA & LP/AB \\
\hline

25--35 & Ubung & Partneruben: gegenseitig Zimmer beschreiben & Dialoge uben & PA & -- \\
\hline

35--40 & Sicherung & Fragen klaren, Tipps fur SL & Fragen stellen & Plenum & Tafel \\
\hline

40--45 & Uber\-leitung & Pause, dann Vorbereitung SL & Aufraumen & -- & -- \\
\hline
\end{longtable}

\vspace{1em}

\textbf{2. Stunde: Stundenleistung (45 Min.)}

\begin{longtable}{|p{1cm}|p{1.5cm}|p{4cm}|p{3cm}|p{2cm}|p{2cm}|}
\hline
\textbf{Zeit} & \textbf{Phase} & \textbf{Lehrerhandeln} & \textbf{Schulerhandeln} & \textbf{SF} & \textbf{Medien} \\
\hline
\endhead

0--5 & Vorbe\-reitung & SL austeilen, Regeln erklaren, Zeit: 25 Min. & Zuhoren, Material bereitlegen & Plenum & SL \\
\hline

5--30 & Stunden\-leistung & Aufsicht, Fragen nur zur Verstandlichkeit & SL bearbeiten & EA & SL \\
\hline

30--35 & Ein\-sammeln & SL einsammeln & Abgeben & -- & -- \\
\hline

35--40 & Reflexion & Kurzes Feedback: ,,Was fiel leicht/schwer?'' & Ruckmeldung geben & Plenum & -- \\
\hline

40--45 & Abschluss & Ausblick nachste Sequenz, HA & Notieren & Plenum & Tafel \\
\hline
\end{longtable}

\textbf{Legende:} EA = Einzelarbeit, PA = Partnerarbeit, SF = Sozialform

% ============================================================================
% 7. ERWARTETE SCHWIERIGKEITEN
% ============================================================================
\section{Erwartete Schwierigkeiten}

\begin{tabular}{|p{5cm}|p{8cm}|}
\hline
\textbf{Schwierigkeit} & \textbf{Reaktion} \\
\hline
Verwechslung Genus bei Mobeln & Merkregel: -a meist feminin, -o meist maskulin \\
\hline
Adjektivanpassung vergessen & Visualisierung: Endungen farbig markieren \\
\hline
\textit{hay} mit Artikel verwechselt & Kontrastubung: \textit{hay una mesa} vs. \textit{la mesa es grande} \\
\hline
Zeitdruck bei SL & Aufgabe 4 als Erweiterung kennzeichnen \\
\hline
LRS-Schuler & Zeitzugabe (+5 Min.), großere Schrift \\
\hline
\end{tabular}

% ============================================================================
% 8. HAUSAUFGABE
% ============================================================================
\section{Hausaufgabe}

\begin{itemize}
    \item Vokabeln Sequenz 3 wiederholen (Vorbereitung Klassenarbeit)
    \item Optional: Lernpfad LP-03d.html bearbeiten
\end{itemize}

% ============================================================================
% 9. ANLAGEN
% ============================================================================
\section{Anlagen}

\begin{enumerate}
    \item Wiederholungs-Arbeitsblatt (AB-03d.pdf)
    \item Musterlosung Wiederholung (ML-03d.pdf)
    \item Stundenleistung (SL-03d.pdf)
    \item Musterlosung SL (SL-03d-ML.pdf)
    \item Lehrerhinweise (LH-03d.pdf)
    \item Lernpfad (LP-03d.html)
\end{enumerate}

% ============================================================================
% 10. DIDAKTIK-SCORE
% ============================================================================
\newpage
\section{Didaktik-Score}

\begin{scorebox}
\textbf{Bewertung der didaktischen Qualitat nach 10 Dimensionen}\\
\textbf{Ziel-Score: $\geq$ 9.0 (Premium-Qualitat)}

\vspace{1em}
\begin{tabular}{|c|l|c|c|p{4.5cm}|}
\hline
\textbf{\#} & \textbf{Dimension} & \textbf{Gew.} & \textbf{Score} & \textbf{Notiz} \\
\hline
1 & Differenzierung & 15\% & 9/10 & [B]/[S]/[E] in LH, Aufg. 4 gestuft \\
\hline
2 & Taxonomie (AFB) & 10\% & 9/10 & AFB I (20\%), II (60\%), III (20\%) \\
\hline
3 & Lernziel-Alignment & 10\% & 10/10 & Alle Sequenzinhalte abgedeckt \\
\hline
4 & Aktivierung & 10\% & 9/10 & Wiederholung aktiv, dann Test \\
\hline
5 & Scaffolding & 10\% & 9/10 & Wortlisten, Beispiele vorhanden \\
\hline
6 & Feedback-Qualitat & 10\% & 9/10 & Musterlosung, Bewertungskriterien \\
\hline
7 & Sprachliche Klarheit & 10\% & 10/10 & A1-Niveau, klare Aufgaben \\
\hline
8 & Motivation \& Relevanz & 10\% & 9/10 & Alltagsthema Haus/Wohnen \\
\hline
9 & Inklusion (UDL) & 10\% & 9/10 & Zeitzugabe, großere Schrift moglich \\
\hline
10 & Praktikabilitat & 5\% & 10/10 & 25 Min. Test realistisch \\
\hline
\multicolumn{3}{|r|}{\textbf{GESAMT:}} & \textbf{9.2/10} & \\
\hline
\end{tabular}

\vspace{1em}
\textbf{Bewertungsschwellen:}
\begin{itemize}
    \item[$\boxtimes$] \textcolor{successgreen}{$\geq$ 9.0: \textbf{FREIGABE} (Premium-Qualitat)}
    \item[$\Box$] \textcolor{warningorange}{8.0--8.9: Iteration empfohlen}
    \item[$\Box$] 6.0--7.9: Uberarbeitung erforderlich
    \item[$\Box$] $<$ 6.0: Grundlegende Neukonzeption
\end{itemize}
\end{scorebox}

\end{document}
